\documentclass{article}
\usepackage[landscape]{geometry}
\usepackage{url}
\usepackage{multicol}
\usepackage{amsmath,amsfonts}
\usepackage{esint}
\usepackage{tikz}
\usepackage{amsmath,amssymb}
\usepackage{graphicx}                % include graphics usage
\usepackage{colortbl,xcolor}
\usepackage{mathtools}

\usepackage{enumitem}
\usetikzlibrary{decorations.pathmorphing}
\makeatletter

\newcommand*\bigcdot{\mathpalette\bigcdot@{.5}}
\newcommand*\bigcdot@[2]{\mathbin{\vcenter{\hbox{\scalebox{#2}{$\m@th#1\bullet$}}}}}
\makeatother

\title{130 Cheat Sheet}
\usepackage[brazilian]{babel}
\usepackage[utf8]{inputenc}

\advance\topmargin-.8in
\advance\textheight3in
\advance\textwidth3in
\advance\oddsidemargin-1.5in
\advance\evensidemargin-1.5in
\parindent0pt
\parskip2pt
\newcommand{\hr}{\centerline{\rule{3.5in}{1pt}}}
%\colorbox[HTML]{e4e4e4}{\makebox[\textwidth-2\fboxsep][l]{texto}

%------------------------------------------
% TEST SECTION ON COPYRIGHT FOOTNOTE
%------------------------------------------


%------------------------------------------
% BODY OF DOCUMENT
%------------------------------------------
\begin{document}

\begin{center}{\huge{\textbf{ESPEC GL Controller Software Ver. 4.0.0 Reference Sheet}}} \\ {\bf{ESPEC North America, INC. $\bullet$ 4141 Central Pkwy $\bullet$ Hudsonville, MI 49426 (USA)}} \\
\end{center}
\begin{multicols*}{3}

\tikzstyle{mybox} = [draw=black, fill=white, very thick,
    rectangle, rounded corners, inner sep=10pt, inner ysep=10pt]
\tikzstyle{fancytitle} =[fill=black, text=white, font=\bfseries]

%-----------------------------------------
% Block or Page Section
%-----------------------------------------

%------------ ESPEC Web Controller ---------------------
\begin{tikzpicture}
\node [mybox] (box){%
\begin{minipage}{0.3\textwidth}
ESPEC Web Controller V3 (EWC V3) is an embedded computer powered by 
Debian GNU/Linux. It utilizes a Web browser to display 
(i.e., host) its user interface (UI) for chamber operation. EWC V3 
can operate as a standalone machine with its touchscreen to control the chamber or 
as part of a network. \\

{\bf{Security:}} Both EWC V3 and Debian continually receive security updates. 
Its root system is configured with read-only access to prevent intruders 
from exploiting the system. Only authorized users from a PC on the same 
network can access and control EWC V3. \\   

{\bf{Stability:}} EWC V3 uses a dual-root (stepladder) technique to update its
firmware. One root system is always in a stable state to ensure an uninterrupted operation. \\

{\bf{Access Privilege \& User Policy:}} EWC V3 is shipped with one user account (called admin)
with complete powers to manage the system. {\bf{admin}} can enforce a user policy 
and administer additional accounts with various privileges. 

\end{minipage}
};
%------------ Title ----------------------------------
\node[fancytitle, right=10pt] at (box.north west) {What is ESPEC Web Controller (EWC)?};
\end{tikzpicture}
%------------ ESPEC Web Controller ---------------------

%------------ Software/Hardware Requirement ------------
\begin{tikzpicture}
\node [mybox] (box){%
    \begin{minipage}{0.3\textwidth}
No software installation is required to use EWC V3; only a 
standard Web browser on the local PC is required  
to access and control EWC V3. \\

{\bf{Web browser:}} Firefox, Chromium, Google Chrome, Safari, 
Opera, MS IE (ver. 10+) or Edge. \\ 

{\bf{OS:}} MS Windows, Linux, Mac OS X, Android, Apple. \\

{\bf{Network Requirement:}} EWC V3 and PC must be on the same network. 
Network is optional for T-series.
    \end{minipage}
};
%------------ Title bar ------------------------------
\node[fancytitle, right=10pt] at (box.north west) {OS \& Software/Hardware Requirement};
\end{tikzpicture}
%------------ Software/Hardware Requirement ------------


%------------ Controller Support ---------------------
\begin{tikzpicture}
\node [mybox] (box){%
\begin{minipage}{0.3\textwidth}
%EWC V3 supports a variety of PLCs: \\ {\bf{ESPEC}} P300, SCP220, ES102;
EWC V3 supports the following PLCs: \\ {\bf{ESPEC}} P300, SCP220, ES102;  
    {\bf{Watlow}} F4/F4T; {\bf{Allen Bradley}} ControlLogix, 
    CompactLogix and Micro8xx series; available for T-series chambers only.
\end{minipage}
};
%------------ Title ----------------------------------
\node[fancytitle, right=10pt] at (box.north west) {PLC Compatibility};
\end{tikzpicture}
%------------ Controller Support ---------------------

%------------ Touchscreen HMI Support ---------------------
\begin{tikzpicture}
\node [mybox] (box){%
\begin{minipage}{0.3\textwidth}
EWC V3 supports a touchscreen or a standard display for operation. 
T-series and custom chambers are shipped with EWC V3 and a touchscreen as their default HMI. 
A touchscreen or a standard display allows EWC V3 to function as a standalone system without 
joining a network.
\end{minipage}
};
%------------ Title bar ----------------------------------
\node[fancytitle, right=10pt] at (box.north west) {Touchscreen or Standard Display};
\end{tikzpicture}
%------------ Touchscreen HMI ---------------------


%------------ Chamber Interface Configuration ---------------------
\begin{tikzpicture}
\node [mybox] (box){%
   \begin{minipage}{0.3\textwidth}
     TCP/IP and Serial interfaces are two different methods used by EWC V3 to
     communicate with chamber/PLC.

     \begin{itemize}
        \item {\bf{TCP/IP:}} T-series chamber (Allen Bradley PLC), Watlow F4T. 
             This protocol offers fast communication between PLC and EWC V3.  
        \item {\bf{Serial:}} ESPEC P300, SCP220, ES102, Watlow F4. 
             Default serial interface for P300 is RS-485; SCP220, ES102 and F4 use RS-232C. 
        \item {\bf{Serial-to-USB:}} A communication interface for a Universal 
             Web Controller kit (U-Web kit). The U-Web kit is available for purchase to operate on 
             an existing chamber. 
     \end{itemize} 
      
   \end{minipage}
};
%------------ Title bar -------------------------------------------
\node[fancytitle, right=10pt] at (box.north west) {Chamber Communication Interface};
\end{tikzpicture}
%------------ Chamber Interface Configuration ---------------------


%------------ Network Configuration & Hostname --------------------
\begin{tikzpicture}
\node [mybox] (box){%
\begin{minipage}{0.3\textwidth}
EWC V3 supports both DHCP and Static networks. 

   \begin{itemize}
      \item {\bf{DHCP:}} EWC V3 automatically joins customer DHCP network, using 
                         a unique IP address.  
      \item {\bf{Static:}} EWC V3 supports various static network configurations. 
             On a small and private (Class C) static network, EWC V3 will use
             its default fallback static IP address {\url{192.168.0.83}} with
             Subnet 255.255.255.0 and Gateway: 192.168.0.1. Refer to User's Manual for detail. 
   \end{itemize} 
On any network, EWC V3 can be accessed via its hostname or IP address. Default 
hostname from factory is: {\url{http://especSN.local/}}, where {\bf{SN}} is 
the serial number of the chamber. 
\end{minipage}
};
%------------ Title Bar --------------------------------------------
\node[fancytitle, right=10pt] at (box.north west) {Network Configuration};
\end{tikzpicture}
%------------ Hostname & Network Configuration ---------------------


%START------------ User Interface --------------------
\begin{tikzpicture}
\node [mybox] (box){%
\begin{minipage}{0.3\textwidth}
EWC V3 offers very intuitive operation menus:
   \begin{itemize}
      \item {\bf{User:}} An authorized user logs in via this menu; user can also reset 
                         their password in this menu.  
      \item {\bf{Overview:}} Chamber current status and operating mode are displayed here. 
                             It is the home page of EWC V3, after a user has logged in.
      \item {\bf{Trend:}} Data points accumulated during chamber operation are displayed as a trend graph. 
                          Data can be downloaded in whole or in portion.
      \item {\bf{History:}} A history of chamber operation modes, its statistics and alarms, is provided. 
      \item {\bf{Constant:}} Constant mode settings of Temperature, Humidity, Refrigeration, Vibration and Outputs can be reset.
      \item {\bf{Program:}} All programmable features of the supported PLCs can be composed into programs to control the chamber. 
                            This menu offers: (1) compose program; (2) open/view program; (3) preview program output prior to execution; 
                            (4) edit and/or overwrite existing program; (5) delete program; (6) rename program; 
                            (7) download or upload programs, and much more.
      \item {\bf{Start Stop:}} This menu allows an authorized operator to control or manage the chamber in any available 
                               operation modes: Standby, Constant and Program.
      \item {\bf{Settings:}} Different configuration settings of EWC V3 and chamber, such as hardware monitor 
                             (for T-series), network, email, user interface, data logging, date/time, 
                             user accounts, macro and controller settings, chamber interface, firmware, 
                             API and Server, can be controlled under this menu.   
      \item {\bf{About:}} EWC V3 firmware version, end-user information and User's Manual give users full control of the system. 
   \end{itemize} 
 
\end{minipage}
};
%------------ Title Bar --------------------------------------------
\node[fancytitle, right=10pt] at (box.north west) {User Interface \& Operation Menus};
\end{tikzpicture}
%END------------ User Interface ---------------------------------------


%START------------ Operation Mode ---------------------------------------
\begin{tikzpicture}
\node [mybox] (box){%
\begin{minipage}{0.3\textwidth}
Overview page showing chamber in program mode with execution time 
and step being executed. \\

\includegraphics[width=\textwidth]{figures/overview-prog-running.PNG}
%\includegraphics[width=\textwidth]{figures/tseries-prog-run-001.PNG}

\end{minipage}
};
%------------ Title Bar --------------------------------------------
\node[fancytitle, right=10pt] at (box.north west) {Overview};
\end{tikzpicture}
%END------------ Operation Mode ---------------------------------------

%------------ Monitor Page -----------------------------------------
\begin{tikzpicture}
\node [mybox] (box){%
\begin{minipage}{0.3\textwidth}
List of chamber operation statistics with alarms: \\

\includegraphics[width=\textwidth]{figures/history-00.PNG}

\end{minipage}
};
%------------ Title Bar --------------------------------------------
\node[fancytitle, right=10pt] at (box.north west) {History};
\end{tikzpicture}
%------------ Monitor Page -----------------------------------------

%START------------ PROGRAM Page -----------------------------------------
\begin{tikzpicture}
\node [mybox] (box){%
\begin{minipage}{0.3\textwidth}
Program can be previewed prior to execution, thus allowing operator
to see its output in detail.  \\

\includegraphics[width=\textwidth]{figures/program-preview-00.PNG}

\end{minipage}
};
%------------ Title Bar --------------------------------------------
\node[fancytitle, right=10pt] at (box.north west) {Program};
\end{tikzpicture}
%------------ Monitor Page -----------------------------------------

%------------ start Stop Menu ------------------------------------------
\begin{tikzpicture}
\node [mybox] (box){%
\begin{minipage}{0.3\textwidth}
Operator can quickly initiate any operating mode: \\

\includegraphics[width=\textwidth]{figures/start-stop-menu-00.PNG}

Operator may step through the program to study the effects of instructions
in a specific step. 

\end{minipage}
};
%------------ Start Stop Header -----------------------------------
\node[fancytitle, right=10pt] at (box.north west) {Start Stop};
\end{tikzpicture}
%------------ Start Stop ------------------------------------------

%------------ Chamber Interface Configuration ---------------------
\begin{tikzpicture}
\node [mybox] (box){%
\begin{minipage}{0.3\textwidth}
Interface configuration allows quick \& easy setup: \\

\includegraphics[width=\textwidth]{figures/setting-chamber-interface-00.PNG}

\end{minipage}
};
%------------ Chamber Interface Header ----------------------------
\node[fancytitle, right=10pt] at (box.north west) {Chamber Interface Configuration};
\end{tikzpicture}
%------------ Chamber Interface Configuration ---------------------

%------------ Network Display ---------------
\begin{tikzpicture}
\node [mybox] (box){%
\begin{minipage}{0.3\textwidth}
Quick access to other EWC V3 on the local network:

\includegraphics[width=\textwidth]{figures/network-list-001a.PNG}

\end{minipage}
};
%------------ Network Display Header ---------------------
\node[fancytitle, right=10pt] at (box.north west) {Network Display Page};
\end{tikzpicture}

%------------ Rest API/TCP Forwarder ---------------------
\begin{tikzpicture}
\node [mybox] (box){%
\begin{minipage}{0.3\textwidth}

{\bf{TCP forwarder}} permits direct access to the controller through EWC V3 
bypassing Web browser interface, thus allowing operator to talk directly 
to chamber/controller via text command (e.g., PuTTY). \\

\includegraphics[width=\textwidth]{figures/restAPI-001.PNG}

{\bf{Rest API}} provides a universal way to interact with EWC V3 via three security modes based 
on different URIs given at each address of EWC V3. 

\end{minipage}
};
%------------ TCP Forwarder ---------------------
\node[fancytitle, right=10pt] at (box.north west) {Rest API/TCP Forwarder};
\end{tikzpicture}

%------------ Firmware Update ---------------
\begin{tikzpicture}
\node [mybox] (box){%
\begin{minipage}{0.3\textwidth}
Each EWC V3 is registered for cloud service support with risk tolerant and
efficient over-the-air updates.  
With customer consent, EWC V3 with access to the Internet will receive firmware update through Mender Cloud service provider.  
\end{minipage}
};
%------------ Firmware Update Header ---------------------
\node[fancytitle, right=10pt] at (box.north west) {Firmware Update \& Cloud Service};
\end{tikzpicture}
%
\columnbreak 

%------------ E-mail Alert ---------------
\begin{tikzpicture}
\node [mybox] (box){%
\begin{minipage}{0.3\textwidth}
To guard against extended damage to the chamber when a certain 
condition arises, an operator can set up an e-mail alert notification. 
Using the Simple Mail Transfer Protocol or SMTP, EWC V3 
can notify operators via e-mail regarding the conditions of the 
chamber(s) such as when alarms went off. In order for this feature to 
function, EWC V3 must have access to the Internet. 
Multiple e-mail addresses can be configured to notify different 
recipients about the conditions of the chamber(s). 

\end{minipage}
};
%------------ E-mail Alert Header ---------------------
\node[fancytitle, right=10pt] at (box.north west) {E-mail Alert};
\end{tikzpicture}
%

%------------ Custom Actions & Macro Editor ---------------
\begin{tikzpicture}
\node [mybox] (box){%
    \begin{minipage}{0.3\textwidth}
Custom actions related to chamber's condition and its operation 
state can be created to monitor and send out an e-mail alert 
when a certain operation state has reached. A feature known 
as {\bf{Macro Editor}} allows operators to configure a set of 
custom actions that will trigger and send out e-mail alerts when 
a specific condition occurs. Here are examples of the available custom actions: 
\begin{itemize}
\item E-mail alert upon program completion
\item E-mail alert when an alarm goes off
\item E-mail alert when chamber changes its operation state 
(e.g., from On to Off, Standby or Constant, etc.)
\item E-mail alert on control loop parameter
\item E-mail alert on output time signals
\item E-mail alert when date/time has changed
\item Upload data logs to an MS Windows SAMBA share server
\item Upload data logs via FTP server
\item Custom trigger (operators' custom scripts)
\item E-mail alert when parameters from external devices are 
written to the Web Cotnroller via Modbus TCP/RTU
\item upload/download custom action configuration
\end{itemize} 

    \end{minipage}
};
%------------ E-mail Alert Header ---------------------
\node[fancytitle, right=10pt] at (box.north west) {Macro Editor \& Custom Actions};
\end{tikzpicture}
%
{\small {
\copyright Reference Sheet, Rev. 1.0: January 2025. 
Features and options on this reference sheet for ESPEC GL Controller, Ver. 4.0.3, 
are subject to change without notice. 
}}
\end{multicols*}
\end{document}
